% Noether P06 Korean producer draft — UNCHECKED
% T11-U149; ED0001 line5401; source 952 LF B / 2E61340138EE5F9B7B190A67EA63B16FF574232150AAC0804B629E5D12C938C7
% Pointer v008 / ED0001: F13E6D896DE6403829FE902609668AB3E3FCA8C3C7FAA07BE5F7A7A72A4C33D8 / D1F06B311F6CBD991DD247D745DD9A72DDE326A20396DF43CFE0C8EDB1593CDB
% Translation only; no review/build/render/approval. Monic irreducible form, uniqueness, transfer, SatzVI denominator collapse, and integrity-basis conclusion remain checker holds.

실제로 이 정의에 따르면 부정원 $x_1\cdots x_\rho$와 따라서 선형형식 $u(x)=u_1x_1+\cdots+u_\rho x_\rho$도 $\Gdom^*_{\rho\rho}$에 대하여 대수적으로 정수적으로 종속된다. 그러므로 근 $z=u(x)$를 갖는 기약 정함수 --- 이를 포함하는 최소 체 $\Kfield^*_{\rho\rho}$의 인볼루션형식 --- 는 최고차항의 계수가 1이고 나머지 계수들은 $\Gdom^*_{\rho\rho}$의 다항식이다. 따라서 이는 동시에 $\Gdom^*_{\rho\rho}$의 인볼루션형식이며, $\Gdom^*_{\rho\rho}$에는 다른 인볼루션형식이 존재할 수 없다. 전이 원리에 따라 $\Gdom_{n\rho}$도 최고차항의 계수가 1인 인볼루션형식을 갖는다. 정리 VI에 따라 $\Gdom_{n\rho}$의 모든 다항식을 그에 대응하는 인볼루션기저로 유리적으로 나타낼 때 분모에는 이 최고차항의 계수만 들어가므로, 인볼루션기저는 정기저가 된다. 즉,
